% 原书第 129—132 页；6.11。
\section{条件分布的算法独立性}
\label{sec:ch06-aic}

第 6.10 节说明，因果结构不仅蕴含统计意义上的（条件）独立关系，也蕴含相对于其他非统计信息度量的独立关系。我们还看到，可以为算法信息表述马尔可夫条件。于是，算法马尔可夫条件最基本的推论，就是 Reichenbach 原理在算法相依关系中的对应版本：只有当两个对象有共同原因，或其中一个影响另一个时，它们才可能在算法意义上相依［Janzing and Schölkopf, 2010］。否则，它们会被空集 d-分离，从而相互独立。类似地，因果上互不相关的 $d$ 个对象 $x_1,\ldots,x_d$，在算法意义上相互独立，即
\begin{equation}
 K(x_1,\ldots,x_d)\overset{+}{=}\sum_{j=1}^dK(x_j).
 \tag{6.29}
\end{equation}
类似于统计信息论中的术语，也可以将左右两侧之差称为\textbf{多元信息（multi-information）}，并将相互独立写为
\begin{equation}
 I(x_1,x_2,\ldots,x_d)\overset{+}{=}0.
 \tag{6.30}
\end{equation}
相互独立还蕴含每个子集的独立性。例如，如果 $x_1,x_2$ 的联合描述，比二者分别描述更短，那么 $x_1,\ldots,x_d$ 的联合描述，也必然比所有 $x_j$ 分别描述更短。因此，式 (6.30) 蕴含
\[
 I(x_1:x_2)\overset{+}{=}0.
\]

如果现在假设，因果图模型中的条件分布\footnote{如前所述，记号 $P_{Y\mid X}$ 是条件分布族 $(P_{Y\mid X=x})_x$ 的简写。}$P_{X_j\mid\PA_j}$ 是“由自然界独立选择的”，就可以推断它们在算法意义上相互独立［Janzing and Schölkopf, 2010；Lemeire and Janzing, 2013］，从而得到原理 4.13 的多元版本。

\begin{principle}[条件分布的算法独立性（AIC）]
在定义 6.21 第 (iii) 项所述的因果贝叶斯网络中，由马尔可夫核描述的因果条件分布，在算法意义上相互独立，即
\begin{equation}
 I(P_{X_1\mid\PA_1},P_{X_2\mid\PA_2},\ldots,P_{X_d\mid\PA_d})
 \overset{+}{=}0,
 \tag{6.31}
\end{equation}
等价地，
\begin{equation}
 K(P_{X_1,\ldots,X_d})\overset{+}{=}\sum_{j=1}^dK(P_{X_j\mid\PA_j}).
 \tag{6.32}
\end{equation}
\end{principle}

注意，不要将原理 6.53 与第 6.10 节的算法马尔可夫条件混淆。后者涉及 $n$ 个单独对象之间的因果关系，不涉及统计抽样；前者则仍然假定 $n$ 个随机变量的传统独立同分布情形，只是额外提出一项推断原理。

与二元情形相同，式 (6.31) 与式 (6.32) 的等价性是直接的，因为描述联合分布，等价于描述所有因果马尔可夫核。换言之，AIC 表示：联合分布的最短描述，由各因果马尔可夫核的分别描述给出。

因果忠实性与 AIC 的思想相近，往往给出类似结论。为讨论相似与不同之处，重新考察例 6.34。参数 $a$ 描述 $P_{Y\mid X}$，参数 $(b,c)$ 描述条件分布 $P_{Z\mid X,Y}$，因此
\begin{equation}
 I(P_{Y\mid X}:P_{Z\mid X,Y})\overset{+}{\ge}I(a:(b,c)).
 \tag{6.33}
\end{equation}
这是因为，只关注两个对象的某些“方面”，不能增加它们之间的算法互信息，例如见 Janzing and Schölkopf［2010，引理 6］。当线性模型的结构系数满足
\begin{equation}
 a\cdot b+c=0
 \tag{6.34}
\end{equation}
时，会出现“非一般的”独立关系 $X\indep Z$。此时 $K(a\mid b,c)\overset{+}{=}0$，因为可以利用长度为 $O(1)$ 的程序，从 $b,c$ 计算出 $a$。因此，
\[
 I(a:(b,c))\overset{+}{=}K(a)-K(a\mid(b,c)^*)\overset{+}{=}K(a).
\]
可见，只要 $K(a)$ 显著大于 $0$，就违反 AIC。对于一个一般的实数 $a$，$K(a)$ 随要求的相对精度提高而按对数增长。于是，AIC 会拒绝相应的因果 DAG，因为式 (6.34) 被视为一种不大可能发生的巧合。

需要解释“只要 $K(a)$ 显著大于 $0$”这句话，因为它反映了 AIC 与忠实性在概念上的不同。设 $b=c$、$a=-1$，则式 (6.34) 成立，但知道 $b,c$ 并不能进一步缩短 $a$ 的描述，因为 $K(a)$ 本来就可忽略。因此，尽管式 (6.34) 看似体现了参数的精细调节，AIC 却并未被违反。沿用 Lemeire and Janzing［2013］的思路，下面说明为何我们认为不拒绝这类调节是 AIC 的特点，而非缺陷。

关键想法是：在给定精度下，结构系数 $\pm1$ 在自然界中出现的频率，远高于 $2.36724\ldots$ 这样的“更一般”数值。例如，花费金额 $S$，会使可用金额 $A$ 减少 $S$，即变化为 $-S$。因此，$S$ 与 $A$ 之间的因果关系由结构系数 $-1$ 描述。\footnote{本例提示，结构系数简单，往往源于我们定义变量的方式，而非“自然界”本身的性质。一般而言，可以追问：我们在多大程度上刻意以能够产生简单因果关系的方式定义变量？}
隐含地，AIC 及上述论证基于一种先验：描述长度较短的取值更可能出现，这与 Solomonoff 的归纳推断理论［Solomonoff, 1964］一致。

AIC 的另一项特点是，它也会拒绝不同路径\emph{近乎}抵消的情形。例如，假设 $a$ 非常接近 $-c/b$。为估计此时的 $I(a:(b,c))$，注意到
\[
 I(a:(b,c))\overset{+}{\ge}I(a:(c/b)),
\]
并使用以下想法：两个相近整数 $n,m$ 的算法互信息，通常约为 $\log\bigl(n/|m-n|\bigr)$，因为知道 $m$ 后，描述 $n$ 约需 $\log|n-m|$ 位，否则约需 $\log n$ 位。进行任意精细的离散化后，可以用整数表示 $a$ 与 $c/b$，进而以 $\log[a/(a+c/b)]$ 粗略估计 $P_{Y\mid X}$ 与 $P_{Z\mid X,Y}$ 之间的算法互信息。
\begin{quote}\small
\textbf{译注：}最后这个粗略估计按原书保留。若将其理解为由两个数的接近程度节省的描述长度，应以差异的绝对值衡量距离；相应地，通常按 $\log|a/(a+c/b)|$ 理解，以避免参数符号使实对数无定义。
\end{quote}
\FloatBarrier
